% Drawing 152/01-10 -- folder 152, batch 01, archive page 1 (PDF page 2).
% Transcription: au-dessous, au crayon : un arbre à trois branches (un trait épais descend d'un point marqué à un gros sommet, d'où partent deux branches vers le bas, celle de gauche épaisse et terminée par un gros point) ; à droite, une flèche montante d'où part un trait ; un « W » de quatre flèches issues de deux pointes basses ; un losange de quatre flèches ; un trait et une boucle en épingle fléchée.
% Kind: three-branched tree, W of arrows, lozenge of arrows, hairpin loop. Method 2 (schematic redraw from the facsimile). Status: ok.
% Notes: Left: the tree with three branches: marked point at the top (it is the lower marked point of 01-9), heavy stroke down to the heavy vertex, a heavy branch to the lower left ending in a heavy point, a thinner branch to the lower right with a point just beyond its end; two faint straight strokes beside it and an isolated point at the right. Then an arrow going up with a stroke leaving its head to the upper right; a W of arrows (two arrows going up from the lower left tip, one short arrow between them, one arrow going up from the lower right tip, a short stroke closing the tip); a lozenge of four arrows (two lower ones going up and outwards, two upper ones going up and inwards); a stroke coming down to a hairpin loop with an arrowhead on its right arm and a short stroke inside it, and a faint curve arriving from the left. Overlaid on the 400 dpi crop. Not drawn: the ink letter I above the W (text). Redrawn 2026-11-12 (issue 128): to check again.
% Coordinates: 1 unit = 1 mm = 1/100 of the crop width (bbox in page fractions: [0.15, 0.29, 0.43, 0.405]).
% Redrawn by Claude (claude-fable-5-2), issue #128, 2026-11-12; not checked by a human; not compiled (no TeX in the session).
% Rights: reproduces part of the work -- behind RENDER_TRANSCRIPTS, never published by itself.
\documentclass[tikz,border=3pt]{standalone}
\usepackage{amsmath,amssymb}
\usetikzlibrary{patterns,arrows.meta}
\tikzset{pencil/.style={black!45}, struck/.style={black!35}, mpt/.style={circle,draw,inner sep=1.3pt}, pt/.style={circle,fill,inner sep=0.9pt}, opt/.style={circle,draw,fill=white,inner sep=1.1pt}}
\begin{document}
\begin{tikzpicture}[x=1mm,y=1mm,line cap=round,line join=round,font=\small]
\node[pt] at (24,-23.5) {};
\draw[very thick] (24,-23.5) -- (23,-35.5);
\node[pt] at (23,-35.5) {};
\draw[thick] (23,-35.5) -- (5.6,-51);
\node[pt] at (5.6,-51) {};
\draw (23,-35.5) -- (33.6,-49);
\node[pt] at (35.2,-51.5) {};
\draw[pencil] (25,-20.5) -- (32,-29.5);
\draw[pencil] (10,-27) -- (5.6,-49);
\node[pt] at (34.5,-25) {};
\draw[-{Stealth[length=1.8mm]}, pencil] (39.6,-25.5) -- (39,-13);
\draw[pencil] (39,-11.6) -- (47.5,-7.5);
\draw[-{Stealth[length=1.8mm]}, pencil] (56.6,-40.2) -- (55.4,-26.2);
\draw[-{Stealth[length=1.8mm]}, pencil] (60.6,-32.6) -- (57.4,-27.8);
\draw[-{Stealth[length=1.8mm]}, pencil] (56.6,-40.2) -- (66.6,-27.4);
\draw[-{Stealth[length=1.8mm]}, pencil] (65.6,-38.6) -- (67.6,-30.4);
\draw[pencil] (63.4,-35.6) -- (65.2,-39.6);
\draw[-{Stealth[length=1.8mm]}, pencil] (76,-17.6) -- (82.6,-8.4);
\draw[-{Stealth[length=1.8mm]}, pencil] (88,-19) -- (84.6,-9.6);
\draw[-{Stealth[length=1.8mm]}, pencil] (79.6,-30) -- (76.6,-19);
\draw[-{Stealth[length=1.8mm]}, pencil] (82.6,-28.4) -- (88.6,-21);
\draw[pencil] (88.6,-33.6) -- (86,-42);
\draw[pencil] plot[smooth, tension=0.55] coordinates {(75.6,-39.6) (81,-40.2) (85.6,-42)};
\draw[-{Stealth[length=1.8mm]}, pencil] plot[smooth, tension=0.55] coordinates {(84,-43.6) (83,-50) (83.4,-53.4) (85.2,-53.8) (86.4,-50) (87.4,-45)};
\draw[pencil] (84.6,-45) -- (84.6,-51);
\end{tikzpicture}
\end{document}
